% ===================================================================
%  TAVOLA N. 5 — La casa e come si costruisce.
%  Traduction 2026 d'apres texto/io, controlee sur texto/fr et
%  texto/en. Consigne : voir texto/it/10-tavola-01.tex.
%
%  Le tableau 5 est un DIALOGUE, et les deux volumes composent le nom
%  du parleur en petites capitales : on garde \textsc{}. Les renvois
%  du plan sont des LETTRES, et l'on ecrit ce que l'ido ecrit, y
%  compris le « (1) » compose pour « (l) ».
%
%  LE TABLEAU 5 ALLEMAND AVAIT MONTRE QUE GUIGNON S'ETAIT TROMPE DE
%  MOT. CELUI-CI MONTRE QU'IL AVAIT LE BON DANS SA PROPRE LISTE.
%
%  La note de l'alinea 25 propose la racine « balk-o » et l'appuie sur
%  six langues : « Balk-o = allemand Balken, anglais joist, francais
%  solive, ITALIEN TRAVICELLO, russe balka, espagnol et portugais
%  viga ». Puis elle definit ce qu'elle veut dire : la solive, « qui
%  n'est ni une trabo (francais poutre) NI UNE PETITE POUTRE ».
%
%  Or « travicello » est exactement une petite poutre : c'est le
%  diminutif de « trave », qui est la poutre, le « trabo » qu'il
%  refuse. Guignon cite donc, dans sa propre liste, un mot BATI SUR
%  L'ANALYSE QU'IL DECLARE FAUSSE. Et il choisit finalement le mot
%  allemand, « Balken », qui designe justement la grosse poutre
%  porteuse — c'est-a-dire l'autre chose qu'il refuse.
%
%     SES DEUX TEMOINS LE CONTREDISENT PAR LES DEUX BOUTS : l'italien
%     dit que la solive est une petite poutre, l'allemand qu'elle est
%     une grosse. Il ecarte l'un et prend l'autre, et il n'obtient ni
%     ce qu'il voulait ni ce qu'il croyait avoir.
%
%  La colonne ecrit « travicello » a l'alinea 25, parce que c'est le
%  mot italien de la chose ; la note, elle, garde « travicello » aussi,
%  parce que c'est le mot que la note cite. Pour une fois les deux
%  coincident, et l'on n'a rien a declarer.
%
%  ET L'IDRAULICO DE L'ALINEA 28 EST LE SEUL PLOMBIER D'EUROPE QUI NE
%  SOIT PAS NOMME PAR LE PLOMB. Le « plumbarius » romain travaillait
%  les tuyaux de plomb, et il a donne le plombier francais, le plumber
%  anglais, le fontanero espagnol mis a part. L'italien l'appelle
%  idraulico, du grec « hydraulikos », qui etait l'orgue a
%  eau. LE MOT EST PLUS RECENT QUE LES AUTRES ET DECRIT UNE PLOMBERIE
%  PLUS RECENTE : il date de l'eau courante sous pression, non des
%  conduites de plomb. Une langue nomme un metier a la date ou elle le
%  regarde, et le tableau 5 allemand disait deja la meme chose de ses
%  corporations.
%
%  Le reste des metiers est latin et transparent : muratore celui qui
%  fait le mur, falegname celui qui travaille le bois, imbianchino
%  celui qui blanchit, vetraio celui du verre, fabbro le faber,
%  scalpellino celui du ciseau, manovale celui de la main. AUCUN N'EST
%  UN NOM DE CORPORATION, tous sont des noms d'action : la ou
%  l'allemand avait grave un statut de ville du XIIIe siecle,
%  l'italien decrit simplement ce que l'homme fait.
% ===================================================================

%%K t05-apar-1 apar
\VUsaut{6.0mm}
\VUtitre{15.0pt}{60}{TAVOLA N. 5}
\VUsaut{1.4mm}
\VUtitre{11.4pt}{0}{La casa e come si costruisce.}
\VUsaut{2.2mm}

%%K t05-tit-1 sub
\VUpk{10.2pt}{40}{Un incontro fortunato.}
\VUsaut{3.0mm}

%%K t05-01-1 p
\textsc{Giovanni}. --- Zio, mi piacerebbe proprio sapere come si costruisce una
\VUgras{casa}.

%%K t05-01-2 p
\textsc{Lo zio} \textsuperscript{(80)}. --- È facile, ragazzo mio. Vieni con me e ti
faccio vedere quella che il signor Bernardo \textsuperscript{(79)} sta facendo
costruire e in cui andrò ad abitare.

%%K t05-01-3 p
\textsc{Giovanni}. --- E chi ha fatto il \VUgras{progetto} \textsuperscript{(76)} di
questa casa?

%%K t05-01-4 p
\textsc{Lo zio}. --- L'\VUgras{architetto} \textsuperscript{(75)}; l'ha disegnato in
modo chiarissimo, è stato approvato, e l'\VUgras{impresario}
\textsuperscript{(77)} ha cominciato i lavori.

%%K t05-01-5 p
\textsc{Giovanni}. --- E chi è l'impresario?

%%K t05-01-6 p
\textsc{Lo zio}. --- Il signor Bianchi, il padre del tuo amico Giacomo. Dirige gli
operai insieme al signor Roux, l'architetto.

%%K t05-01-7 p
Ma guarda! Ecco il signor Bianchi che arriva proprio al momento giusto con il suo
\VUgras{metro} \textsuperscript{(78)}, e il signor Roux con il suo progetto.

%%K t05-01-8 p
Buongiorno, signori. Come state?

%%K t05-01-9 p
\textsc{Il signor Roux}. --- Benissimo, grazie. Buongiorno, Giovanni; com'è cresciuto
questo ragazzo!

%%K t05-01-10 p
\textsc{Lo zio}. --- Eh sì, è già un ometto. È molto serio; tutto quello che vede vuole
che gli sia spiegato. Oggi vorrebbe sapere come si costruisce una casa.

%%K t05-tit-2 sub
\VUpk{10.2pt}{40}{Gli operai.}
\VUsaut{3.0mm}

%%K t05-01-11 p
\textsc{Il signor Bianchi}. --- Allora vieni con me, che te lo spiego subito. Mi
piacciono i bambini che vogliono imparare.

%%K t05-01-12 p
Vedi quell'operaio con la \VUgras{pala} \textsuperscript{(82)} in mano: è uno
\VUgras{sterratore} \textsuperscript{(81)}. Sta scavando una \VUgras{trincea}
\textsuperscript{(46)} per posare la conduttura dell'acqua. È lui che ha scavato anche
il terreno su cui sta la casa, perché il muratore potesse alzare i quattro
\VUgras{muri}. La parte di quei muri che sta nel terreno si chiama le
\VUgras{fondamenta} \textsuperscript{(2)}. Lo spazio racchiuso dalle fondamenta forma
la \VUgras{cantina} \textsuperscript{(3)}. Questa cantina ha una finestrella che si
chiama \VUgras{sfiatatoio} \textsuperscript{(5)}, una \VUgras{porta}
\textsuperscript{(4)} e una scala.

%%K t05-01-13 p
Il \VUgras{muratore} \textsuperscript{(108)} ha poi continuato ad alzare i muri,
lasciando le aperture per le \VUgras{porte} \textsuperscript{(8)} e le
\VUgras{finestre} \textsuperscript{(17)}.

%%K t05-01-14 p
\textsc{Giovanni}. --- Ma io non lo vedo, questo muratore!

%%K t05-01-15 p
\textsc{Il signor Bianchi}. --- Con la casa ha finito. È là dalla \VUgras{stalla}
\textsuperscript{(54)}, su per la sua \VUgras{impalcatura} \textsuperscript{(110)}, e
sta tirando su il muro del \VUgras{lavatoio} \textsuperscript{(56)}.

%%K t05-01-16 p
Ha una cazzuola, una cassa da malta e un \VUgras{filo a piombo}
\textsuperscript{(109)}. Ma questi nomi non te li ricorderai.

%%K t05-01-17 p
\textsc{Giovanni}. --- Sì che me li ricordo, perché chiederò a mio zio una cazzuola
piccola e una cassa da malta, e il filo a piombo me lo faccio io con un pezzo di
spago.

%%K t05-tit-3 sub
\VUsaut{3.0mm}
\VUpk{10.2pt}{40}{I materiali.}
\VUsaut{3.0mm}

%%K t05-01-18 p
\textsc{Lo zio}. --- Benissimo. Ma dimmi un po', Giovanni, che cosa serve per
costruire una casa?

%%K t05-01-19 p
\textsc{Giovanni}. --- Servono \VUgras{pietra da taglio} \textsuperscript{(59)} e
\VUgras{malta} \textsuperscript{(65)}.

%%K t05-01-20 p
\textsc{Lo zio}. --- Esatto. Guarda quell'uomo con il cappellone, sul suo
\VUgras{carro} \textsuperscript{(136)}. È un \VUgras{carrettiere}
\textsuperscript{(135)}. Tutti i giorni porta qui \VUgras{blocchi di pietra}
\textsuperscript{(60)}. Scarica il carro nel cortile, e gli \VUgras{scalpellini}
\textsuperscript{(83)} squadrano i blocchi e li segano con la loro
\VUgras{sega da pietra} \textsuperscript{(85)}. Un \VUgras{manovale}
\textsuperscript{(146)} li porta al muratore, che li mette in opera.

%%K t05-01-21 p
Intanto il \VUgras{maltaio} \textsuperscript{(87)} prepara la malta. Gli servono
\VUgras{sabbia} \textsuperscript{(62)}, \VUgras{calce} \textsuperscript{(63)} e
\VUgras{acqua} \textsuperscript{(64)}. La calce è accanto a lui in un \VUgras{sacco}
\textsuperscript{(71)}; un \VUgras{garzone muratore} \textsuperscript{(111)} gli porta
la sabbia in una \VUgras{carriola} \textsuperscript{(112)}, e l'\VUgras{apprendista}
\textsuperscript{(113)} è alla pompa \textsuperscript{(58)}. Sta riempiendo un
\VUgras{secchio} \textsuperscript{(114)}. La malta sarà pronta fra poco. Il
\VUgras{portamalta} \textsuperscript{(107)} la prende e la porta su per la scala fino
all'impalcatura, dove un muratore sta finendo quel lato della casa.

%%K t05-01-22 p
E adesso, come si chiama l'operaio che fa il \VUgras{tetto} \textsuperscript{(31)}?

%%K t05-01-23 p
\textsc{Giovanni}. --- È il \VUgras{conciatetti} \textsuperscript{(118)}.

%%K t05-tit-4 sub
\VUsaut{3.0mm}
\VUpk{10.2pt}{40}{Il tetto della casa.}
\VUsaut{3.0mm}

%%K t05-01-24 p
\textsc{Il signor Bianchi}. --- Sì, ma prima viene il \VUgras{carpentiere}
\textsuperscript{(115)}.

%%K t05-01-25 p
Il carpentiere fa l'\VUgras{orditura del tetto} \textsuperscript{(35)}. Unisce le
\VUgras{travi} \textsuperscript{(37)} con \VUgras{caviglie} \textsuperscript{(117)}.
Quegli operai lassù, con i larghi calzoni di velluto, sono carpentieri. Uno tiene una
piccola trave e l'altro taglia una caviglia con la sua \VUgras{ascia}
\textsuperscript{(116)}. Quello vicino al \VUgras{camino} \textsuperscript{(41)} è il
conciatetti. Inchioda le assicelle sui \VUgras{travicelli} (*) e le
\VUgras{lastre di ardesia} \textsuperscript{(66)} sulle assicelle. A volte mette le
tegole.

%%K t05-noto-1 noto
\VUnotes{143.2mm}{%
(*) \VUgras{Balk-o} = tedesco \textit{Balken}, inglese \textit{joist}, francese
\textit{solive}, italiano \textit{travicello}, russo \textit{balka}, spagnolo e
portoghese \textit{viga}. Una radice tecnica non ancora adottata, ma qui proposta per
rendere il senso della parola francese \textit{solive}, che non è né una
\textit{trabo} (francese \textit{poutre}) né una piccola trave. È un legno squadrato
che porta un pavimento e un soffitto.%
}

%%K t05-01-26 p
Il tetto della stalla è di \VUgras{tegole} \textsuperscript{(55)}, quello della casa è
di \VUgras{ardesia} \textsuperscript{(32)}, e quello della veranda è di
\VUgras{zinco} \textsuperscript{(24)}.

%%K t05-01-27 p
\textsc{Giovanni}. --- I tetti di zinco li fa anche il conciatetti?

%%K t05-01-28 p
\textsc{Il signor Bianchi}. --- No, quello è il \VUgras{lattoniere}
\textsuperscript{(121)} o l'\VUgras{idraulico}, quello che sta mettendo un
\VUgras{abbaino} \textsuperscript{(38)} nel tetto della casa. È lui che mette anche le
\VUgras{grondaie} \textsuperscript{(43)} e i \VUgras{pluviali}
\textsuperscript{(42)}. Salda lo zinco e la latta. Ed è stato lui a fissare il
\VUgras{parafulmine} \textsuperscript{(40)} e la \VUgras{banderuola}
\textsuperscript{(39)} sul \VUgras{frontone} \textsuperscript{(36)}.

%%K t05-01-29 p
\textsc{Giovanni}. --- E allora la casa è finita?

%%K t05-tit-5 sub
\VUsaut{3.0mm}
\VUpk{10.2pt}{40}{Gli attrezzi.}
\VUsaut{3.0mm}

%%K t05-01-30 p
\textsc{Il signor Bianchi}. --- Non del tutto. Lo \VUgras{stuccatore}
\textsuperscript{(99)} costruisce con i \VUgras{mattoni} \textsuperscript{(61)} le
pareti che la dividono nelle varie stanze. Mette il \VUgras{gesso}
\textsuperscript{(70)} sui \VUgras{soffitti} \textsuperscript{(26)} e sui muri.
Adesso sta lavorando al soffitto della \VUgras{veranda} \textsuperscript{(33)}. Come
il muratore ha una \VUgras{cazzuola} \textsuperscript{(100)} e una
\VUgras{cassa da malta} \textsuperscript{(101)}.

%%K t05-01-31 p
Poi il falegname fa le porte, le finestre, i \VUgras{pavimenti di legno}
\textsuperscript{(16)} e le \VUgras{persiane} \textsuperscript{(19)}.

%%K t05-01-32 p
Adesso sta lavorando in cortile al suo \VUgras{banco} \textsuperscript{(90)}. Ha una
\VUgras{sega} \textsuperscript{(94)} per tagliare il \VUgras{legno}
\textsuperscript{(69)}, una \VUgras{pialla} \textsuperscript{(91)}, un
\VUgras{morsetto} \textsuperscript{(92)}, uno \VUgras{scalpello}
\textsuperscript{(94 \textit{bis})}, un \VUgras{compasso} \textsuperscript{(95)} e un
\VUgras{mazzuolo} di legno \textsuperscript{(95 \textit{bis})}.

%%K t05-01-33 p
\textsc{Giovanni}. --- Oh, ma fare il falegname è ancora più bello che fare il
muratore. Zio, mi compri un banco, una sega e una pialla? Voglio fare una cuccia per
Medoro.

%%K t05-01-34 p
\textsc{Lo zio}. --- Sì, ragazzo mio.

%%K t05-01-35 p
\textsc{Giovanni}. --- Grazie! E chi è quell'uomo lì davanti al portone?

%%K t05-tit-6 sub
\VUsaut{3.0mm}
\VUpk{10.2pt}{40}{La pittura.}
\VUsaut{3.0mm}

%%K t05-01-36 p
\textsc{Il signor Bianchi}. --- È un \VUgras{imbianchino} \textsuperscript{(102)}. È in
piedi sulla sua scala con un barattolo di \VUgras{vernice} \textsuperscript{(103)} e il
suo \VUgras{pennello} \textsuperscript{(104)}. Passa la vernice sul legno del portone
perché duri più a lungo.

%%K t05-01-37 p
È lui che mette anche la carta da parati nelle stanze.

%%K t05-01-38 p
Il \VUgras{tappezziere} \textsuperscript{(107)}, su una \VUgras{scala}
\textsuperscript{(106)}, fuori sul \VUgras{balcone} \textsuperscript{(28)}, sta
montando una \VUgras{tenda da sole} \textsuperscript{(29)}.

%%K t05-01-39 p
L'operaio vicino alla finestra grande è il \VUgras{vetraio} \textsuperscript{(96)}.
Fissa i \VUgras{vetri} \textsuperscript{(18)} con lo \VUgras{stucco}
\textsuperscript{(73)}. Quell'altro operaio, in ginocchio davanti alla porta della
cantina, è un \VUgras{fabbro ferraio} \textsuperscript{(97)}. Sta montando una
\VUgras{serratura} \textsuperscript{(6)}. La sua \VUgras{lima}
\textsuperscript{(98)} è accanto a lui.

%%K t05-01-40 p
\textsc{Giovanni}. --- E l'uomo che batte un pezzo di ferro con il suo
\VUgras{martello} \textsuperscript{(84)}, è anche lui un fabbro ferraio?

%%K t05-01-41 p
\textsc{Il signor Bianchi}. --- Non proprio --- è un \VUgras{maniscalco}
\textsuperscript{(125)}. Sta forgiando \VUgras{ferramenta} \textsuperscript{(67)} per
l'\VUgras{elettricista} \textsuperscript{(124)}, che è sul tetto della
\VUgras{rimessa} \textsuperscript{(51)}. Ha una \VUgras{forgia}
\textsuperscript{(126)}, un'\VUgras{incudine} \textsuperscript{(127)} e una
\VUgras{tenaglia} \textsuperscript{(128)}.

%%K t05-01-42 p
L'elettricista sta fissando il \VUgras{filo di rame} \textsuperscript{(44)} del
telefono.

%%K t05-tit-7 sub
\VUpk{10.2pt}{40}{Il giardinaggio.}
\VUsaut{3.0mm}

%%K t05-01-43 p
\textsc{Lo zio}. --- In fondo al \VUgras{cortile} \textsuperscript{(45)}, vicino alla
\VUgras{cancellata} \textsuperscript{(47)} e al \VUgras{portone carraio}
\textsuperscript{(48)}, vedi il \VUgras{giardiniere} \textsuperscript{(129)}. Sta
preparando il piccolo \VUgras{giardino} \textsuperscript{(49)}. Ha vangato la terra con
la sua \VUgras{vanga} \textsuperscript{(131)}, sta seminando e rastrellando con il
\VUgras{rastrello} \textsuperscript{(130)}. Poi, con il suo \VUgras{annaffiatoio}
\textsuperscript{(132)}, annaffierà le \VUgras{aiuole} \textsuperscript{(150)} e le
piante cresceranno.

%%K t05-01-44 p
Il \VUgras{palafreniere} \textsuperscript{(133)}, che ha appena governato e nutrito il
cavallo, sta pulendo la \VUgras{carrozza} \textsuperscript{(52)} con una spugna, una
\VUgras{pompa a mano} \textsuperscript{(134)} e un secchio d'acqua. Poi la metterà
nella rimessa e chiuderà la porta con il pesante \VUgras{catenaccio}
\textsuperscript{(53)}.

%%K t05-01-45 p
Ecco, spero che adesso sarai contento; di spiegazioni ne hai avute abbastanza.

%%K t05-01-46 p
\textsc{Giovanni}. --- Sì, zio, ma mi piacerebbe vedere anche la casa di dentro.

%%K t05-01-47 p
\textsc{Lo zio}. --- Quello sarà per domani. Il signor Roux ti spiegherà il progetto e
ti dirà come si chiama ogni stanza.

%%K t05-tit-8 sub
\VUsaut{3.0mm}
\VUpk{10.2pt}{40}{La distribuzione della casa.}
\VUsaut{2.4mm}

%%K t05-apar-2 apar
{\centering\textit{(Vedi il progetto.)}\par}
\VUsaut{2.4mm}

%%K t05-01-48 p
\textsc{L'architetto}. --- Vieni, Giovanni, ti porto a vedere le varie parti della
casa.

%%K t05-01-49 p
Entriamo al \VUgras{pianterreno} \textsuperscript{(7)}. Ecco il pulsante del
\VUgras{campanello} \textsuperscript{(11)}. Saliamo i tre \VUgras{gradini}
\textsuperscript{(14)} della \VUgras{scalinata} \textsuperscript{(9)}, varchiamo la
\VUgras{soglia} \textsuperscript{(10)} del \VUgras{portone} \textsuperscript{(8)}, ed
eccoci ai piedi della \VUgras{scala} \textsuperscript{(13)}.

%%K t05-01-50 p
\textsc{Giovanni}. --- Questo è il corridoio.

%%K t05-01-51 p
\textsc{L'architetto}. --- No, questo è l'\VUgras{ingresso} \textsuperscript{(12)}. Il
\VUgras{corridoio} \textsuperscript{(a)} è in fondo. A destra ci sono il
\VUgras{salotto} \textsuperscript{(b)} e la \VUgras{sala da pranzo}
\textsuperscript{(c)}; di fronte alla sala da pranzo la \VUgras{cucina}
\textsuperscript{(d)} e la \VUgras{dispensa} \textsuperscript{(e)}; poi, sul davanti,
lo \VUgras{studio} \textsuperscript{(f)}. La \VUgras{veranda}
\textsuperscript{(23)}, con la sua \VUgras{porta a vetri} \textsuperscript{(25)}, è
accanto al salotto.

%%K t05-01-52 p
Adesso saliamo le scale. Tieniti al \VUgras{corrimano} \textsuperscript{(15)} e conta i
gradini. Sono quattordici. Ecco il \VUgras{pianerottolo} \textsuperscript{(m)}: siamo
al primo \VUgras{piano} \textsuperscript{(27)}.

%%K t05-tit-9 sub
\VUsaut{3.0mm}
\VUpk{10.2pt}{40}{Il primo piano.}
\VUsaut{3.0mm}

%%K t05-01-53 p
Di fronte alle scale ci sono il \VUgras{gabinetto} \textsuperscript{(20)} e lo
\VUgras{spogliatoio} \textsuperscript{(j)}. A destra una bella \VUgras{camera da
letto} \textsuperscript{(g)} con due finestre sulla \VUgras{facciata}
\textsuperscript{(1)} della casa. A sinistra il \VUgras{bagno} \textsuperscript{(1)} e
la \VUgras{camera degli ospiti} \textsuperscript{(i)} con una grande
\VUgras{alcova} \textsuperscript{(h)}. Accanto alla camera da letto c'è la
\VUgras{camera dei bambini} \textsuperscript{(k)}. Dà su un'ampia \VUgras{terrazza}
\textsuperscript{(21)}. Sotto questa terrazza c'è un altro gabinetto
\textsuperscript{(20)}, e al muro c'è la \VUgras{campana} \textsuperscript{(22)} che
suona per il pranzo.

%%K t05-01-54 p
\textsc{Giovanni}. --- Saliamo al \VUgras{secondo piano}.

%%K t05-01-55 p
\textsc{L'architetto}. --- Un secondo piano non c'è. Sotto il tetto ci sono solo le
\VUgras{soffitte} \textsuperscript{(33)} e il solaio \textsuperscript{(34)}. Il giro è
finito --- possiamo scendere.

%%K t05-01-56 p
\textsc{Giovanni}. --- Grazie, signore, è stato davvero interessante. Vado a raccontare
a mio padre tutto quello che ho visto.
\VUsaut{4.0mm}
\VUfilet{22mm}
